\chapter{Concluding remarks\label{chap9}}


Although the problems addressed in this thesis are part of very active topics in physics and engineering, we believe that many important questions remain unanswered. \\
In a significant part of this thesis we used simplified models, such as the coupled dipole model. These models have the advantage of explaining complex physical phenomena by simple principles. However, they are not as accurate as other methods when the goal is to reproduce experiments. Thus, this thesis is directed to the explanation and prediction of phenomena and not their accurate reproduction.\\
%In this thesis we tackle some open problems, presenting now the main conclusions.\\
In this chapter we present the main conclusions of our work.\\
The first part of the thesis, composed by chapters (\ref{chap3}), (\ref{chap4}) and (\ref{chap5}), is devoted to the study of light scattering phenomena in magnetoplasmonic systems, while the second part, composed by chapters (\ref{chap6}), (\ref{chap7}) and (\ref{chap8}), we study light emission statistics in random correlated systems.\\

In chapter (\ref{chap3}) we present a study of the influence of the magneto-optical activity over a nanostructure. We consider a dimer system formed by two nanodisks, one plasmonic and other magnetoplasmonic, where the electromagnetic coupling is controlled via distance between them. The system is illuminated and a magnetic field is applied perpendicularly to the surface of the nanodisks, \textit{i.e.}, in Kerr configuration.
 As consequence of the electromagnetic excitation and the application of the magnetic field, an induced net electric dipole along the perpendicular direction (perpendicular to the incident light and its polarisation direction, and to the applied magnetic field) in the magneto-optic active nanodisk is generated. Due to the interference between both nanodisks, a magneto-optical activity in the plasmonic nanodisk is induced. Also the interference of the fields from the nanodisks give rise to an electromagnetically induced magneto-optical transparency in these systems. This leads to a Fano-like spectral dependence of the MO activity.\\
In chapter (\ref{chap4}) we continue the study of dimer systems but this time presenting a detailed analysis of the interaction effects when the plasmonic or magneto-plasmonic nature of the nanodisk components is changed. We conclude that, for specific configurations, the magneto-optical response can be dominated by the induced magneto-optical activity of the purely plasmonic nanodisk. The magneto-optical activity of a system with only one of the nanodisks containing material with intrinsic magneto-optics can be even larger than that of a system composed of two magnetoplasmonic nanodisks.\\
%We tackle the problem by considering a simple dipole model, where the global phenomenology of the system is captured. The results are supported by experimental and full numerical calculations.\\
In chapter (\ref{chap5}) we present a study of the properties of an emitter in two different configurations when an external magnetic field is applied: in the presence of a single magnetoplasmonic nanoparticle and inside of a cavity formed by two magnetoplasmonic nanoparticles. Both systems, for a specific distance between the particles and the emitter, present a minimum in the radiative decay rate.
At this minimum, the presence of an external magnetic field leads to a large modification of the radiated field patterns, due to the opening of new radiation channels. The work undertaken in this chapter serves as a first approach to problems where the control of light emission is done by magneto-optical effects. \\
The study of simple magneto-optical systems like those described above, does not only explain the physics behind these specific problems, but also encourages the study of other systems. The manipulation of the geometry of the particles, shapes and sizes, separately or together, can result in a diverse and rich phenomenology. Magneto-Chiral, disordered or correlated magnetoplasmonic systems are some of the unexplored problems with great potential.

In chapter (\ref{chap6}) we study the statistical decay rate dependence of a point emitter in a system of nanoparticles statistically distributed according to a 2D lattice-gas model at different temperatures. The temperature, in this problem, is the responsible for the control of the correlations between vacancies. When the ordering temperature is far from criticality, the correlation length is small, and the decay rate distribution has the typical long-tailed shape where events with large Purcell factors are rare.
%Is shown that for short-range correlations, the Purcell factors present non-Gaussian long-tailed statistics where events with large Purcell factors are extremely rare.
For temperatures near the critical point, where the correlation length tends to infinity, the statistics evolves towards a bimodal distribution with two well-defined peaks, one at high enhancement factors, while another peak remains at values corresponding to free space. %We show that diluted systems, where vacancies are distributed in a long-range manner, enclose clusters of dipoles of all sizes, allowing for the existence of many configurations with optical modes confined around the source.\\
For diluted systems we show that it is possible that many configurations with optical modes confined around the source exist.\\
In chapter (\ref{chap7}) we study the self-diffusion of a strongly confined Lennard-Jones system. For clusters with a few hundred of particles, the system present a dynamic phase switching between solid-like and liquid-like phases. The monitoring of the self-diffusion coefficient in the solid-like and liquid-like phase in the phase-switching region, indicates a difference of a factor of three between both phases (with a larger value in the liquid-like phase). Unexpectedly, the radial distribution function in the phase-switching region is essentially indistinguishable between both phases.\\% In this way, a possible way to identify phases could be made by the analysis of the internal energy evolution or the self-diffusion characteristics of the system, because structural information does not provide a relevant signature of phase switching in this case.\\
In chapter (\ref{chap8}) we propose a method based on emission decay rate experiments to identify the different phases. Although light scattering experiments do not discriminate the system in the switching regime, single emitter decay-rate statistics experiments (with low diffusive emitters) may allow the identification of phases. This behaviour can be attributed to the difference in the radial distribution function between scatterers and the emitter which, in turn, might also be attributed to differences in the self-diffusion of scatterers between both phases. This work proposes a new way of monitor subtle thermodynamical behaviours of complex systems at sizes comparable with the wavelength of the light source employed in the experiment.\\
%The understanding of light scattering and emission through complex systems is as interesting as complex and difficult. Despite the efforts, much remains to be done. Introduction of optical forces in complex systems or experimental implementation of two and three-dimensional systems are some of the unexplored paths.
The understanding of light scattering and emission through complex systems is as interesting as complex and difficult. Despite the efforts, much remains to be done. Questions such as the importance of optical forces in complex systems or the introduction of nonreciprocal effects in disordered systems are actual open problems.
